% Options for packages loaded elsewhere
\PassOptionsToPackage{unicode}{hyperref}
\PassOptionsToPackage{hyphens}{url}
\PassOptionsToPackage{dvipsnames,svgnames,x11names}{xcolor}
%
\documentclass[
  11pt,
  a4paper,
]{ltjsarticle}
\usepackage{amsmath,amssymb}
\usepackage{iftex}
\ifPDFTeX
  \usepackage[T1]{fontenc}
  \usepackage[utf8]{inputenc}
  \usepackage{textcomp} % provide euro and other symbols
\else % if luatex or xetex
  \usepackage{unicode-math} % this also loads fontspec
  \defaultfontfeatures{Scale=MatchLowercase}
  \defaultfontfeatures[\rmfamily]{Ligatures=TeX,Scale=1}
\fi
\usepackage{lmodern}
\ifPDFTeX\else
  % xetex/luatex font selection
\fi
% Use upquote if available, for straight quotes in verbatim environments
\IfFileExists{upquote.sty}{\usepackage{upquote}}{}
\IfFileExists{microtype.sty}{% use microtype if available
  \usepackage[]{microtype}
  \UseMicrotypeSet[protrusion]{basicmath} % disable protrusion for tt fonts
}{}
\makeatletter
\@ifundefined{KOMAClassName}{% if non-KOMA class
  \IfFileExists{parskip.sty}{%
    \usepackage{parskip}
  }{% else
    \setlength{\parindent}{0pt}
    \setlength{\parskip}{6pt plus 2pt minus 1pt}}
}{% if KOMA class
  \KOMAoptions{parskip=half}}
\makeatother
\usepackage{xcolor}
\usepackage[margin=24mm]{geometry}
\setlength{\emergencystretch}{3em} % prevent overfull lines
\providecommand{\tightlist}{%
  \setlength{\itemsep}{0pt}\setlength{\parskip}{0pt}}
\setcounter{secnumdepth}{5}
\ifLuaTeX
  \usepackage{selnolig}  % disable illegal ligatures
\fi
\IfFileExists{bookmark.sty}{\usepackage{bookmark}}{\usepackage{hyperref}}
\IfFileExists{xurl.sty}{\usepackage{xurl}}{} % add URL line breaks if available
\urlstyle{same}
\hypersetup{
  colorlinks=true,
  linkcolor={blue},
  filecolor={Maroon},
  citecolor={Blue},
  urlcolor={blue},
  pdfcreator={LaTeX via pandoc}}

\author{}
\date{}

\begin{document}

\hypertarget{ux884cux5217ux5206ux89e3ux306eux5168ux4f53ux50cf}{%
\section{行列分解の全体像}\label{ux884cux5217ux5206ux89e3ux306eux5168ux4f53ux50cf}}

行列分解の目的は、複雑な線形写像を、意味の分かる単純な写像の積へ分けることです。

代表例は

\begin{itemize}
\tightlist
\item
  LU: 消去法の構造を分離する
\item
  QR: 直交基底と座標係数へ分離する
\item
  固有値分解: 固有方向と伸縮率へ分離する
\item
  SVD: 入力側の直交基底・伸縮・出力側の直交基底へ分離する
\end{itemize}

です。

\hypertarget{ux5206ux89e3ux306fux540cux3058ux3082ux306eux3092ux5225ux306eux5ea7ux6a19ux3067ux898bux308bux6280ux8853}{%
\subsection{分解は「同じものを別の座標で見る」技術}\label{ux5206ux89e3ux306fux540cux3058ux3082ux306eux3092ux5225ux306eux5ea7ux6a19ux3067ux898bux308bux6280ux8853}}

行列 \texttt{A}
を直接眺めても、その本質が見えにくいことがあります。分解は、\texttt{A}
がどの方向に何をしているのかを露出させます。

例えば SVD

\[
A=U\Sigma V^\top
\]

は、

\begin{enumerate}
\def\labelenumi{\arabic{enumi}.}
\tightlist
\item
  \texttt{V\^{}T} で入力を右特異ベクトル基底へ変換する
\item
  \texttt{Σ} で各軸を特異値だけ伸縮する
\item
  \texttt{U} で出力空間へ向きを合わせる
\end{enumerate}

という3段階です。

\hypertarget{ux5206ux89e3ux3092ux9078ux3076ux57faux6e96}{%
\subsection{分解を選ぶ基準}\label{ux5206ux89e3ux3092ux9078ux3076ux57faux6e96}}

「最良の分解」が一つあるわけではありません。

\begin{itemize}
\tightlist
\item
  方程式を何度も解く → LU
\item
  最小二乗を安定に解く → QR
\item
  正方行列の反復挙動を見る → 固有値分解
\item
  ランク・条件数・擬似逆・低ランク近似を統一して見る → SVD
\end{itemize}

と目的に応じて選びます。

\end{document}
